% Part: methods
% Chapter: induction
% Section: introduction

\documentclass[../../../include/open-logic-section]{subfiles}

\begin{document}

\olfileid{mth}{ind}{int}

\olsection{పరిచయం}

ఆగమనం ఒక ముఖ్యమైన నిరూపణ పద్ధతి. తర్కం, సిద్ధాంతాత్మక కంప్యూటర్ శాస్త్రం, గణితశాస్త్రం దాదాపు అన్ని రంగాలలో అది వివిధ రూపాల్లో ఉపయోగపడుతుంది. తర్కంలోని అనేక ఫలితాలను నిరూపించడానికి ఇది అవసరం.

ఆగమనాన్ని తరచుగా నిగమనంతో పోల్చి, ప్రత్యేక సందర్భాల నుంచి సాధారణ విషయానికి వెళ్లే నిగమనంగా వర్ణిస్తారు. ఉదాహరణకు మనం అనేక పచ్చని ఎమరాల్డ్ రత్నాలను చూసి, పచ్చగా లేని ఏ ఎమరాల్డ్‌నూ చూడకపోతే, అన్ని ఎమరాల్డ్‌లు పచ్చనివని తేల్చవచ్చు. ఇది ఆగమన నిగమనం: అనేక ప్రత్యేక సందర్భాల నుంచి (ఈ ఎమరాల్డ్ పచ్చనిది, ఆ ఎమరాల్డ్ పచ్చనిది మొదలైనవి) ఒక సాధారణ వాదనకు (అన్ని ఎమరాల్డ్‌లు పచ్చనివి) వెళ్తుంది. \emph{గణిత} ఆగమనం కూడా సాధారణ వాదనకు చేరే నిగమనమే; కానీ ఈ ``సరళ ఆగమనం'' కంటే చాలా భిన్నమైనది.

స్థూలంగా చెప్పాలంటే గణితంలో ఆగమన నిరూపణ ఒక రకానికి చెందిన అన్ని గణిత వస్తువులకు ఒక లక్షణం ఉందని నిర్ధారిస్తుంది. సరళమైన సందర్భంలో అవి సహజ సంఖ్యలు. అప్పుడు ఆగమన నిరూపణ అన్ని సహజ సంఖ్యలకు ఒక లక్షణం ఉందని స్థాపించడానికి ఈ రెండింటిని చూపుతుంది:
\begin{enumerate}
    \item $0$కు ఆ లక్షణం ఉంది; మరియు
    \item ఒక సంఖ్య~$k$కు ఆ లక్షణం ఉంటే,~$k+1$కూ ఉంటుంది.
\end{enumerate}
సహజ సంఖ్యలపై ఆగమనాన్ని, సంఖ్యలను కేటాయించగల గణిత వస్తువుల గురించి సాధారణ వాదనలను నిరూపించడానికి కూడా తరచుగా ఉపయోగించవచ్చు. ఉదాహరణకు ప్రతి పరిమిత సమితిలో పరిమిత సంఖ్యలో, అంటే~$n$ మూలకాలు ఉంటాయి. ప్రతి సంఖ్య~$n$కు ``పరిమాణం~$n$ గల అన్ని పరిమిత సమితులు \dots'' అనే లక్షణం ఉందని ఆగమనంతో చూపగలిగితే, అన్ని పరిమిత సమితుల గురించీ ఒక విషయాన్ని చూపినట్లవుతుంది.

ఆగమనాన్ని \emph{ఆగమనాత్మకంగా నిర్వచించిన} గణిత వస్తువులకు కూడా సాధారణీకరించవచ్చు. ఉదాహరణకు మొదటిస్థాయి తర్క వ్యక్తీకరణల వంటి ఆకారిక భాష వ్యక్తీకరణలు ఆగమనాత్మకంగా నిర్వచించబడతాయి. \emph{నిర్మాణాత్మక ఆగమనం} అటువంటి అన్ని వ్యక్తీకరణల గురించి ఫలితాలను నిరూపించే పద్ధతి. ముఖ్యంగా నిర్మాణాత్మక ఆగమనం తర్కంలో చాలా ఉపయోగకరమూ, విస్తృతంగా వాడేదీ.

\end{document}
